\chapter{Troubleshooting}\label{ch:troubleshooting}

\begin{longtable}{@{}P{0.30\linewidth} P{0.30\linewidth} P{0.32\linewidth}@{}}
\caption{Symptoms, causes and remedies.}\label{tab:troubleshooting}\\
\toprule
Symptom & Cause & Remedy \\
\midrule
\endfirsthead
\toprule
Symptom & Cause & Remedy \\
\midrule
\endhead
\val{source the OpenFOAM bashrc first} & No OpenFOAM environment. & \Cref{ch:setup}. \\
Negative cell volumes. & The surface breaks a rule of \cref{tab:surface-rules}, or a level is too coarse for a curvature. & Do not use the mesh. Correct the surface or the levels. \\
A wall patch without layers. & No entry in \ent{patchBoundaryLayers}. & \Cref{sec:layers}. \\
Incorrect domain, or a number with a decimal comma. & \cmd{awk} in a locale with a decimal comma. & \Cref{sec:ts-locale}. \\
\cmd{rhoCentralFoam} fails in the first iterations. & Start from a uniform free stream. & \Cref{sec:supersonic-start}. \\
Courant number of order $10^7$ in the log. & Normal with \ent{localEuler}. & \Cref{sec:ts-courant}. \\
A run with \ent{--init} behaves as a start from rest. & \cmd{mapFields} failed. & \Cref{sec:ts-mapfields}. \\
\cmd{taskset} shows all CPUs for a process; many migrations. & The binding does not work. & \Cref{sec:binding}. \\
Wall time much larger than CPU time; \ent{suspendedSeconds} $> 0$. & The computer suspended. & \Cref{sec:prepare-computer}. The coefficients are correct; the cost is not. \\
\ent{profiles} has more than one value; low \ent{GHzEffective}. & The CPU profile changed during the run. & \Cref{sec:prepare-computer}. \\
No \ent{cycles} in \file{run.info}. & \ent{perf_event_paranoid} too high. & \Cref{sec:prepare-computer}. \\
The computer swaps or the run stops for memory. & The mesh is too large for the memory. & Scale the memory from a coarser mesh (\cref{tab:requirements}). \\
\cmd{pkill -f} stops your own shell. & The command line contains the pattern. & Stop the process by its number. \\
A sweep reads wrong lines. & The sweep file was edited while it ran. & Stop, edit, start again. \\
A script does unexpected things after an edit. & \cmd{sh} reads a script while it runs it. & Do not edit a script that runs. \\
A dictionary does not parse near its banner. & \texttt{*/} in the banner comment closes the comment. & Remove \texttt{*/}. \\
A script continues after a failed application. & \cmd{runApplication} returns 0 on \val{FOAM FATAL ERROR}. & Read each log, as \file{Allmesh} and \file{Allrun} do. \\
\bottomrule
\end{longtable}

\section{mapFields on the same mesh}\label{sec:ts-mapfields}

\cmd{mapFields} can stop with a segmentation fault when the source and the target are
the same mesh. Between different meshes it works. \file{Allrun} stops the run if the
log of \cmd{mapFields} or \cmd{changeDictionary} does not end with \val{End}; a check
for \val{FOAM FATAL} only does not detect a segmentation fault. On the same mesh, a copy
of the fields of the source run is an alternative that does not interpolate.

\section{The Courant number with a local time step}\label{sec:ts-courant}

With \ent{localEuler}, \cmd{rhoCentralFoam} calculates the Courant number with the
global \ent{deltaT}, so the value in the log is meaningless. Read
\val{Flow time scale min/max} (\cref{sec:lts}).

\section{Decimal comma}\label{sec:ts-locale}

In a locale with a decimal comma, \cmd{awk} reads and writes numbers with a comma, and
OpenFOAM reads them incorrectly. Each script that calculates with \cmd{awk} sets
\cmd{LC_ALL=C}. Do the same in a new script.

\section{Rules for a change to the scripts}

\begin{itemize}
\item Measure before you propose a change; a change without a number is not complete.
\item Do not put a Python script in a bash here-document: a \texttt{\textbackslash n}
  in an f-string becomes a real new line and breaks the file.
\item After a change, run one short case and examine \file{run.info}.
\item Change this guide in the same commit.
\end{itemize}
